\subsection{自由代数的复习}

设 X 为集合。回忆在 X 上构造的自由原群 M(X) 的作法（《代数》，第 I 章，第 7 节，第 1 号）。对整数 \(n \geqslant 1\) 用归纳法，我们如下定义诸集合 \(\mathbf{X}_n\)：写 \(\mathbf{X}_1 = \mathbf{X}\)，并取 \(\mathbf{X}_n\) 为诸集合 \(\mathbf{X}_p \times \mathbf{X}_{n-p}\)（其中 \(p = 1, 2, \ldots, n-1\)）的和集；若 X 有限，则每个 \(\mathbf{X}_n\) 也有限。族 \((\mathbf{X}_n)_{n \geqslant 1}\) 的和集记作 M(X)；每个集合 \(\mathbf{X}_n\)（特别地 X）都等同于 M(X) 的一个子集。 设 w 与 \(w'\) 属于 M(X)；设 p 与 q 为使 \(w \in \mathbf{X}_p\) 与 \(w' \in \mathbf{X}_q\) 的整数，并设 \(n = p + q\)；有序对 \(\langle w, w' \rangle\) 在 \(\mathbf{X}_p \times \mathbf{X}_{n-p}\) 到 \(\mathbf{X}_n\) 的典范单射下的像记作 \(w.w'\)，称为 w 与 \(w'\) 的乘积。 X 到原群 M 的每个映射都能以唯一方式扩张为 M(X) 到 M 的原群同态。

设 \(w\) 属于 \(\mathrm{M}(\mathbf{X})\)；唯一整数 \(n\)（使 \(w \in \mathbf{X}_n\) 的）称为 \(w\) 的长度，记作 \(l(w)\)。于是 \(l(w.w') = l(w) + l(w')\) 对 \(w, w'\) ∈ \(\mathrm{M}(\mathbf{X})\) 成立。 集合 \(\mathbf{X}\) 是 \(\mathrm{M}(\mathbf{X})\) 中由长度为 1 的元素组成的子集。每个元素 \(w\)（长度 \(\geqslant 2\) 的）都能唯一地写成 \(w = w'.w''\) 的形式。

环 K 上系数的原群 M(X) 的代数记作 \(\operatorname{Lib}(X)\)，当需要指明环 K 时记作 \(\operatorname{Lib}_{K}(X)\)。集合 M(X) 是 K-模 \(\operatorname{Lib}(X)\) 的基，因此 X 将等同于 \(\operatorname{Lib}(X)\) 的一个子集。若 A 是代数，X 到 A 的每个映射都可唯一地扩张为 \(\operatorname{Lib}(X)\) 到 A 的同态（《代数》，第 III 章，第 2 节，第 7 号，命题 7）。

